\scanpage{44}
où
\[
\rho_i:C_n\longrightarrow\underline O_S \tag*{(77)}
\]
est l'hom de $\underline O_S$-Algèbres défini par
\[
\rho_i(T)=t_i, \tag*{(78)}
\]
et alors la formule à prouver est \unsure{cas|un|un cas} particulier
de la suivante, bien connue
\[
 \unsure{\mathcal N|\mathscr N|N}_{C(t_1,\ldots,t_n)/\mathcal O_S}(\deleted{\unsure{f}}w)=\prod_{1\leq i\leq n}\rho_i(w)\quad ]] \tag*{(79)}
\]

Revenons \deleted{\unsure{au théo|au théorème}} au cas $n=2$, et à la correspondance
entre structures d'Algèbres quadratiques
sur $B$, admettant $e_B$ comme él\textsuperscript{t} unité, et
\deleted{structures \inserted{de} \unsure{quad|quadr}} formes quadratiques $q$ sur $B$ satisfaisant
$q(e_B)=1$. Explicitons la structure sur
$\Sigma^2(F)$ de fibré \deleted{vectoriel}\inserted{affine} sur $2F$, de \deleted{base}
module des translations $(L^{\otimes2})_{2F}$. Il
faut d'abord expliciter
\[
\underline b(q)\in\Gamma(2F)\subset\Gamma(\check B=V) \tag*{(80)}
\]
pour $q$ comme dessus. \deleted{\unsure{On fera|On fixe}}

\uline{Prop.} Soit $\varphi$ la forme bilinéaire associée à
$q$, et considérons la forme
\[
\left|\begin{aligned}
\underline b(q)&:B\longrightarrow\underline O_S\\
\underline b(q)(x)&=\varphi(e_B,x)\deleted{\unsure{+}}\;(=\mathrm{tr}_{B/\mathcal O_S}\,x)
\end{aligned}\right. \tag*{(81)}
\]
On a $\underline b(q)(e_B)=\varphi(e_B,e_B)=2q(e_B)=2$, donc
$\underline b(q)$ est une section de $2F\subset\check{\mathbb V}(\check B)$, qui n'est
% archive sheet number: 43
